\section{Conclusion and Future Work}
\label{sec:conclusion}

This article began from an observation about where the economics of software construction are heading. As agentic, ``spec-driven'' lifecycles such as AWS AI-DLC~\cite{awsaidlc}, GitHub Spec Kit~\cite{githubspeckit}, and Kiro~\cite{awskiro} mature, the implementation increasingly becomes a disposable build output that an agent can regenerate on demand, and the durable, scarce artefact becomes the specification of intent. Yet the specifications that drive these emerging lifecycles share three structural defects. The specification is natural language (NL), so the specification-to-code ``compilation'' performed by the agent is unverifiable; the loop is open, in that code is generated and then inspected by eye with no machine gate; and where verification is present at all it is circular, because the same model that writes the code also writes the tests or oracle against which the code is judged~\cite{liu2023evalplus,liu2026oracle}. The lifecycle is being reorganised around the specification as the source of truth, but the specification it is being organised around cannot bear that weight.

\subsection{Summary of the contribution}
\label{subsec:conclusion-summary}

The contribution of this article is a reference architecture that supplies the missing verification spine: a machine-checkable formal-specification contract layer together with a verification loop. The pivot of the architecture is the formal contract, which mediates between the ambiguous human world above it and the precise machine world below it. Above the pivot lies \emph{elaboration}, the human-confirmed translation of NL intent into a formal contract; below it lie \emph{synthesis}, in which an agent produces code from the contract, and \emph{verification}, in which OpenJML extended static checking~\cite{cok2014openjml} discharged by the Z3 SMT solver~\cite{demoura2008z3} certifies the code against the contract over all inputs rather than over a sampled test suite. The organising idea is the dual role of the specification: the same contract is simultaneously the input the agent builds against and the oracle its output is verified against, and heuristic inference removes the spec-authoring bootstrap cost that historically made this dual role impractical.

Around this pivot the article defines a set of mechanisms intended to make the architecture honest rather than merely plausible. The \emph{ratification gap}---that the human used NL precisely because they will not author formal text, yet the contract must be human-owned to be a trustworthy oracle---is closed by ratifying behaviour rather than logic: the contract is rendered as concrete, checkable examples that become pinned oracles, and a Socratic elaboration process surfaces only the boundary decisions (null, empty, overflow, ordering, duplicates, error-versus-exception) the agent had to resolve. An \emph{adversarial elaboration} agent, the elaboration-layer dual of CEGIS~\cite{solarlezama2008sketching}, searches for scenarios consistent with the NL intent but violating the proposed contract, detecting spec weakening and vacuity. An \emph{independent-authority constraint} forbids the agent that writes the code from authoring the contract it is verified against, dissolving the circularity that the reliability literature documents~\cite{liu2023evalplus,zhang2025hallucination,liu2026oracle}; \emph{provenance- and assurance-tagged obligations} make that constraint auditable and the ``verified'' claim precise; and \emph{tiered assurance} degrades a full proof to a bounded check to a test, recording the assurance level rather than silently dropping intent when verification does not terminate. For brownfield codebases---most real software---the contract is recovered rather than authored: a characterisation phase triangulates code, issue tickets, and documentation into a provenance-tagged candidate contract whose disagreements become candidate bugs and the opening agenda of the human dialogue. Crucially, this article reframes the contribution not as a measured improvement but as a development-process capability: it restores a closed, counterexample-driven loop, with an independently owned machine-checkable contract as the source of truth, to a lifecycle that had abandoned it.

\subsection{The instantiation}
\label{subsec:conclusion-instantiation}

The architecture is not merely described but instantiated. The proof-of-concept reuses the author's existing validated assets as trusted components---the JML-Inferrer for draft contracts, the custom OpenJML fork with Z3 as the oracle that never grades itself, the \texttt{AnnotationToJMLConverter}, and the Dockerized verification pipeline---so that the only new code is an orchestration layer and LLM calls for elicitation, contract drafting, and synthesis. The instantiation is exercised on Apache Commons Lang, a genuine three-stream dataset that couples source with the public Apache JIRA \texttt{LANG-\#\#\#\#} project and rich Javadoc, allowing real code-versus-intent discrepancies to be surfaced. This grounds the architecture in the assets and results of the surrounding PhD programme: inferred specifications improve LLM unit-test generation~\cite{edmonds_article1}; specifications can be embedded into compiled artefacts and OpenAPI documents at full roundtrip fidelity~\cite{edmonds_article2}; heuristic inference admits compositional refinement across method boundaries~\cite{edmonds_article3}; and the inference pipeline integrates into CI/CD with diff-only analysis and fail-open semantics~\cite{edmonds_article4}.

\subsection{Limitations}
\label{subsec:conclusion-limitations}

The article is candid about what it does not claim. Inferred specifications skew weak, describing what code does rather than what it should, so ratification remains real labor that the architecture mitigates---but does not eliminate---through inherited default contracts and API-surface prioritisation. Inferred specifications are sometimes over-strong and wrong, generating false candidate bugs that the human-in-the-loop must adjudicate, which is why every recovered obligation is framed as a candidate rather than a verdict. The translation of tickets and documentation into formal clauses without human confirmation is imperfect, and linking those artefacts to code is noisy, so the brownfield results are credible only on a curated module. Verification scalability and decidability limits are real and known: Z3 times out on multiplicative reasoning, array theory, and rich preconditions, which is precisely the contingency tiered assurance exists to handle. Above all, a wrong specification verified against is worse than no specification, which is why the independent-authority constraint and human ratification are load-bearing rather than ornamental. Verifying code against a code-derived specification is characterisation, not a proof of correctness; correctness relative to intent arises only from the intent streams and human ratification.

\subsection{Future work}
\label{subsec:conclusion-future}

Several lines of work follow directly. The first is to complete the proof-of-concept build beyond the de-risked early milestones: milestone M2 generates contracts from NL by few-shot prompting over the author's own verified JML corpus, with the inferrer supplying brownfield drafts, and milestone M3 implements conversational ratification through the Socratic and example-pinning loop. The second is empirical validation on external production codebases beyond the demonstration corpus, to test whether multi-source triangulation and the synthesis--verification loop generalise past a curated module. The third is an industrial action-research case study, conducted with appropriate human research ethics committee (HREC) approval, in which practitioners use the pipeline on their own codebase so that the ratification labor, the candidate-bug yield, and the trust dynamics of retroactive-then-prospective adoption can be observed in situ rather than assumed. The fourth is integration with the CI/CD machinery of the programme's prior work~\cite{edmonds_article4}, so that compositional change propagation becomes a first-class lifecycle event: because contracts compose, a callee-spec change recomputes caller proof obligations and flags which call sites break without running them, making internal refactors and dependency or library version bumps the same operation and turning ``which of my call sites does this upgrade break behaviourally?'' into a routine, machine-answerable query~\cite{edmonds_article3,jayasuriya2024understanding,raemaekers2017semver}.

These threads converge on a single thesis. In a lifecycle where an agent can regenerate the implementation cheaply, the specification is no longer a precursor to the work but the work itself---and it earns that status only by playing both roles at once. As the precise input the agent builds against and the formal oracle its output is verified against, the dual-role specification is what makes an AI-native development process verifiable rather than merely productive, and it is the artefact around which such a process must be organised.
